\chapter{Modes of Convergence}

In the previous chapter, we established the Central Limit Theorem:
when independent random variables with finite variance are summed and standardized by
$\sqrt n$, the resulting distribution approaches the standard normal bell curve
$\Normal(0,1)$. Comparing the Law of Large Numbers and the Central Limit Theorem highlights an important distinction. The
standardized sum does not freeze onto a single number; from trial to trial it
keeps fluctuating with standard deviation~$1$. It is the \emph{shape of
its distribution} that converges, not the numerical values themselves. By contrast, when
we studied the Law of Large Numbers, the running sample average $\bar X_n$ did
not keep fluctuating -- it approached the true mean $\mu$ and
stayed there.

In calculus, there is only one standard way for a sequence of numbers to converge: the
terms $a_n$ get and stay arbitrarily close to a limit $a$. In probability,
the objects that converge are \emph{random variables} -- measurable functions
defined on a sample space, each inducing a probability distribution.

When we say a sequence $X_1, X_2, X_3, \dots$ ``converges to'' some limit $X$,
we have to specify \emph{in what sense} the sequence approaches the limit.
Should the distributions line up? Should the actual realized values line up?
Should they line up on every possible outcome, or just with high probability, or
in mean? Each question leads to a distinct mathematical notion of convergence, and
they form a strict hierarchy: some imply others, and confusing
them can lead to incorrect proofs.

By the end of this chapter, you will understand four principal modes of convergence
-- \emph{almost sure}, \emph{in probability}, \emph{in $L^p$ (in mean)},
and \emph{in distribution}. You will see which implies which, learn how to
apply key operational tools (the continuous mapping theorem and Slutsky's theorem),
derive the delta method for nonlinear transformations of asymptotically normal estimators,
and see how the Law of Large Numbers and the Central Limit Theorem fit into this unified framework.

\chapterline

\section{Four questions, four answers}
\label{ch22:four}

Start with a single concrete story, and ask the four questions one at a time.

\begin{intuitionbox}
You flip a fair coin over and over. After $n$ flips let $\bar X_n$ be the fraction of
heads so far. Everyone has the intuition that $\bar X_n \to 1/2$. Good. But pin down
what that sentence means. There is not one infinite coin-flip experiment sitting in
front of you; there is a whole \emph{universe} $\Omega$ of possible infinite flip
sequences $\omega$. For each $\omega$ you get a whole numerical sequence
$\bar X_1(\omega), \bar X_2(\omega), \dots$. ``$\bar X_n \to 1/2$'' is a claim about
infinitely many sequences at once. \emph{Which} sequences, and \emph{how}, is the
substance.
\end{intuitionbox}

Here are the four ways to make ``$X_n$ gets close to $X$'' precise, from the most
literal to the most relaxed.

\textbf{1. Almost surely (the strongest, most literal).} Look at each outcome
$\omega$. You get an ordinary numerical sequence $X_n(\omega)$. Demand that this
sequence converges to $X(\omega)$ in the ordinary calculus sense -- for
\emph{essentially every} $\omega$. ``Essentially every'' means the set of bad
$\omega$ where it fails has probability zero. This is pointwise convergence of
functions, with a probability-zero pardon.

\textbf{2. In probability.} Forget tracking individual $\omega$'s for all time.
Instead freeze a tolerance $\eps > 0$ and ask: at stage $n$, how likely is $X_n$ to be
farther than $\eps$ from $X$? Demand that this probability shrinks to $0$ as $n$ grows.
The bad set is allowed to \emph{move around} from one $n$ to the next; it just has to
keep shrinking.

\textbf{3. In $L^p$ (in mean of order $p$).} Measure closeness by an average size of
the error: $\E\abs{X_n - X}^p$. Demand that this number goes to $0$. For $p=2$ this
is the familiar mean-squared error of statistics and signal processing.

\textbf{4. In distribution (the weakest, most generous).} Stop asking the $X_n$ to be
near $X$ \emph{as numbers} at all. Only ask that their \emph{distributions} settle
down: the probability of landing in any reasonable region converges to the right
limiting probability. The random variables need not even live on the same probability
space.

\begin{happeningbox}
The deep reason there are four modes and not one: a random variable is two things glued
together -- a \emph{function} $\omega \mapsto X(\omega)$ assigning a numerical value
to each underlying outcome, and a \emph{distribution} (the law that tells us how much
probability lands in each interval of the real line). Modes 1--3 compare the
\emph{functions} (with different rulers: pointwise, in-probability, in-average). Mode 4
throws the underlying sample space away and compares only the \emph{distributions}.
That is why convergence in distribution is so much weaker, and why it alone does not
require the variables to live on the same probability space. Keep these two sides of a
random variable in mind and the whole hierarchy becomes natural instead of memorized.
\end{happeningbox}

\section{Almost sure convergence}
\label{ch22:as}

\begin{definition}[Almost sure convergence]
\label{ch22:def-as}
A sequence of random variables $X_n$ \emph{converges almost surely} to $X$, written
$X_n \asto X$, if
\[
  \P\!\left(\set{\omega : \lim_{n\to\infty} X_n(\omega) = X(\omega)}\right) = 1.
\]
Equivalently, the set of outcomes on which the numerical sequence $X_n(\omega)$ fails
to converge to $X(\omega)$ has probability $0$.
\end{definition}

The phrase ``almost surely'' (a.s.) -- sometimes ``with probability one'' or
``almost everywhere'' in the measure-theory literature -- is doing real work. We are
\emph{not} promising convergence for every single $\omega$. We are tolerating a
negligible exceptional set.

\begin{example}[A typical bad set is genuinely nonempty]
\label{ch22:ex-bad-set}
In the coin-flipping story, the all-heads sequence $\omega^\star = HHHH\cdots$ is a
perfectly legitimate outcome, and along it $\bar X_n(\omega^\star) = 1 \not\to 1/2$.
So convergence really does fail somewhere. But $\set{\omega^\star}$ has probability
$0$ (indeed the set of \emph{all} sequences whose head-frequency misbehaves has
probability $0$, which is exactly the Strong Law of Large Numbers). Almost sure
convergence disregards such measure-zero exceptional sets.
\end{example}

\begin{whybox}
Why allow an exceptional set of measure zero?
In uncountable sample spaces, individual outcomes have probability zero, as we saw in
Example~\ref{ch22:ex-bad-set}. If we demanded convergence at literally every $\omega \in \Omega$,
even the Strong Law of Large Numbers would fail because of outcomes like $HHHH\cdots$.
Since sets of probability zero never change probabilities or expectations --
$\P(A \cup N) = \P(A)$ whenever $\P(N) = 0$ -- ignoring a measure-zero exceptional set
throws away nothing that probability can observe.
\end{whybox}

\subsection{Reformulating a.s. convergence with $\eps$-tubes}
\label{ch22:as-tube}

There is a second, more operational way to write almost sure convergence that makes
its strength visible and sets up the sharp contrast with convergence in probability.
Saying that an ordinary numerical sequence $X_n(\omega)$ converges to $X(\omega)$ means:
for every tolerance $\eps>0$, the values eventually enter the band
$(X(\omega)-\eps,\, X(\omega)+\eps)$ \emph{and stay there forever}.

Translate ``eventually enters and stays'' into set-theoretic probability:
\[
  X_n \asto X
  \quad\Longleftrightarrow\quad
  \text{for every } \eps>0,\quad
  \P\!\left(\bigcap_{N\ge1}\bigcup_{n\ge N} \set{\abs{X_n - X} \ge \eps}\right) = 0.
\]
Deconstruct that double-set operation inside the probability:
\begin{itemize}
  \item The inner union $\bigcup_{n\ge N} \set{\abs{X_n - X} \ge \eps}$ is the event
        that \emph{at least one} violation of size $\eps$ occurs at or after stage $N$.
  \item The outer intersection $\bigcap_{N\ge1}$ demands that violations keep happening
        past \emph{every} horizon $N$ -- in other words, that deviations of size $\eps$
        occur \emph{infinitely often}.
\end{itemize}

Demanding this probability to be zero means: along almost every trajectory $\omega$,
violations happen only finitely many times. Past some random finite stopping time
$N(\omega)$, the sequence is locked permanently inside the $\eps$-tube
(Figure~\ref{ch22:fig-as-tube}).

\begin{figure}[H]
\centering
\begin{tikzpicture}[scale=0.88]
  \fill[sourceblue!8] (0.8,-0.7) rectangle (14.5,0.7);
  \fill[targetgreen!15] (6.5,-0.7) rectangle (14.5,0.7);

  \draw[->,gray!60] (0.5,0) -- (14.8,0) node[right,black] {\footnotesize $n$};
  \draw[->,gray!60] (0.8,-1.6) -- (0.8,1.7) node[above,black] {\footnotesize $X_n(\omega)$};

  \draw[dashed,sourceblue!80!black,thick] (0.8,0.7) -- (14.5,0.7)
    node[right,font=\scriptsize,sourceblue!80!black] {$X(\omega) + \eps$};
  \draw[dashed,sourceblue!80!black,thick] (0.8,-0.7) -- (14.5,-0.7)
    node[right,font=\scriptsize,sourceblue!80!black] {$X(\omega) - \eps$};
  \node[left,font=\scriptsize,gray!70!black] at (0.8,0) {$X(\omega)$};

  \draw[thick,targetgreen!80!black,dash dot] (7,0.06) -- (7,1.4);
  \draw[thick,targetgreen!80!black,dash dot] (7,-0.8) -- (7,-0.24);
  \node[above,font=\footnotesize,targetgreen!80!black] at (7,1.4) {trapping threshold $N(\omega) = 7$};
  \draw[->,targetgreen!80!black,shorten >=2pt] (7,1.35) -- (7,0.75);

  \coordinate (p1) at (1,1.30);
  \coordinate (p2) at (2,-1.20);
  \coordinate (p3) at (3,0.95);
  \coordinate (p4) at (4,-0.85);
  \coordinate (p5) at (5,0.45);
  \coordinate (p6) at (6,0.85);
  \coordinate (p7) at (7,-0.40);
  \coordinate (p8) at (8,0.32);
  \coordinate (p9) at (9,-0.22);
  \coordinate (p10) at (10,0.18);
  \coordinate (p11) at (11,-0.10);
  \coordinate (p12) at (12,0.08);
  \coordinate (p13) at (13,-0.04);
  \coordinate (p14) at (14,0.02);

  \draw[thick,sourceblue!70!black] 
    (p1) -- (p2) -- (p3) -- (p4) -- (p5) -- (p6) -- (p7) -- (p8) -- (p9) -- (p10) -- (p11) -- (p12) -- (p13) -- (p14);

  \foreach \pt in {p1,p2,p3,p4,p6} {
    \fill[trapred] (\pt) circle (2.2pt);
  }
  \fill[sourceblue!80!black] (p5) circle (2.2pt);
  \foreach \pt in {p7,p8,p9,p10,p11,p12,p13,p14} {
    \fill[targetgreen!80!black] (\pt) circle (2.2pt);
  }

  \node[font=\scriptsize,trapred,above right] at (1,1.30) {escapes};
  \node[font=\scriptsize,targetgreen!80!black,align=center] at (10.5,1.1) 
    {\textbf{Permanently trapped:}\\ $\abs{X_n(\omega) - X(\omega)} < \eps$ for all $n \ge N(\omega)$};

  \foreach \x in {1,2,3,4,5,6,7,8,9,10,11,12,13,14} {
    \draw (\x,0.06) -- (\x,-0.06);
  }
  \node[below,font=\scriptsize,fill=white,inner sep=1pt] at (7,-0.08) {$7$};
\end{tikzpicture}
\caption{Almost sure convergence visualized as pathwise lock-in. For any chosen tolerance
$\eps > 0$, an individual sample path $\omega$ may wander outside the $\pm\eps$ tube at
early times (red points). But almost sure convergence guarantees that along almost every
path, there arrives a finite stopping threshold $N(\omega)$ after which the sequence
remains trapped inside the tube forever (green points). The set of paths that escape
infinitely often has probability zero.}
\label{ch22:fig-as-tube}
\end{figure}

\begin{formalbox}
A clean sufficient condition, by far the most useful in practice, comes from the
first Borel--Cantelli lemma: if for every $\eps>0$
\[
  \sum_{n=1}^{\infty} \P\!\left(\abs{X_n - X} \ge \eps\right) < \infty,
\]
then $X_n \asto X$. The intuition is immediate: if the sum of deviation probabilities
is finite, the \emph{expected total number} of deviations over all $n \ge 1$ is
finite ($\E[\sum_n \ind_{\{\abs{X_n-X}\ge\eps\}}] < \infty$). A quantity whose expectation
is finite cannot equal infinity, so deviations can occur only finitely many times along
almost every path. This is the standard technique behind most strong-law proofs:
bound the tail probabilities using Chebyshev or Chernoff inequalities, verify that
the bounds are summable, and almost sure convergence follows at once.
\end{formalbox}

\section{Convergence in probability}
\label{ch22:prob}

\begin{definition}[Convergence in probability]
\label{ch22:def-prob}
$X_n$ \emph{converges in probability} to $X$, written $X_n \pto X$, if for every
$\eps > 0$,
\[
  \lim_{n\to\infty} \P\!\left(\abs{X_n - X} \ge \eps\right) = 0.
\]
\end{definition}

Compare this carefully with the almost-sure definition. There the ``$\eps$-failure'' event was
required to occur only \emph{finitely often along each path}. Here we only require
that, \emph{at each fixed time} $n$, the failure is unlikely -- and that the
probability of a deviation shrinks with $n$. We do not examine entire paths across all time.
The set of exceptional outcomes may vary with $n$; what matters is that its total probability mass tends to zero.

\begin{intuitionbox}
A geometric distinction between the two modes:
\begin{itemize}
\item \textbf{Almost sure convergence:} follows a fixed outcome $\omega$ through all time $n \ge 1$
and demands that the numerical sequence $X_n(\omega)$ stays within an $\eps$-tube around $X(\omega)$ for all sufficiently large $n$ (except on a set of paths of probability zero).
\item \textbf{Convergence in probability:} examines the whole sample space at each fixed index $n$
and demands that the measure of the set where $\abs{X_n - X} \ge \eps$ tends to zero as $n \to \infty$.
\end{itemize}
At each fixed stage $n$, the probability of a large error is small, even if the specific outcomes where large errors occur change from one step to the next.
\end{intuitionbox}

\subsection{The typewriter sequence: in probability but not almost surely}
\label{ch22:typewriter}

The cleanest demonstration that these are truly different notions is the famous
\emph{typewriter} (or ``sliding bump'') sequence. We do it in full, because once you
see it you never again confuse the two modes.

Let $\Omega = [0,1]$ with the uniform probability (length = probability). We build a
sequence of indicator variables by sweeping a window of shrinking width across
$[0,1]$, over and over, like a typewriter carriage returning faster each pass.

\begin{itemize}
\item Pass 1 (width $1$): $X_1 = \ind_{[0,1]}$.
\item Pass 2 (width $1/2$): $X_2 = \ind_{[0,1/2]}$,\quad $X_3 = \ind_{[1/2,1]}$.
\item Pass 3 (width $1/3$): $X_4 = \ind_{[0,1/3]}$,\quad $X_5 = \ind_{[1/3,2/3]}$,\quad
      $X_6 = \ind_{[2/3,1]}$.
\item Pass $k$ (width $1/k$): $k$ indicators of consecutive length-$1/k$ blocks.
\end{itemize}

\begin{figure}[H]
\centering
\begin{tikzpicture}[scale=1]
  \draw[->,gray!70] (0,0) -- (0,2.8);
  % Pass 1
  \begin{scope}[yshift=2.0cm]
    \draw[gray!60] (0,0) -- (6,0);
    \fill[sourceblue!60] (0,0) rectangle (6,0.36);
    \node[left,font=\footnotesize] at (-0.2,0.18) {Pass $1$ ($k=1$)};
    \node[font=\footnotesize,white] at (3,0.18) {$X_1 = \ind_{[0,1]}$};
  \end{scope}
  % Pass 2
  \begin{scope}[yshift=1.2cm]
    \draw[gray!60] (0,0) -- (6,0);
    \fill[sourceblue!60] (0,0) rectangle (3,0.36);
    \fill[trapred!55] (3,0) rectangle (6,0.36);
    \node[left,font=\footnotesize] at (-0.2,0.18) {Pass $2$ ($k=2$)};
    \node[font=\footnotesize,white] at (1.5,0.18) {$X_2 = \ind_{[0,1/2]}$};
    \node[font=\footnotesize,white] at (4.5,0.18) {$X_3 = \ind_{[1/2,1]}$};
  \end{scope}
  % Pass 3
  \begin{scope}[yshift=0.4cm]
    \draw[gray!60] (0,0) -- (6,0);
    \fill[sourceblue!60] (0,0) rectangle (2,0.36);
    \fill[targetgreen!60] (2,0) rectangle (4,0.36);
    \fill[trapred!55] (4,0) rectangle (6,0.36);
    \node[left,font=\footnotesize] at (-0.2,0.18) {Pass $3$ ($k=3$)};
    \node[font=\footnotesize,white] at (1,0.18) {$X_4$};
    \node[font=\footnotesize,white] at (3,0.18) {$X_5$};
    \node[font=\footnotesize,white] at (5,0.18) {$X_6$};
  \end{scope}
  % Base axis with ticks
  \draw[thick] (0,0) -- (6,0);
  \foreach \x/\lbl in {0/{0}, 2/{1/3}, 3/{1/2}, 4/{2/3}, 6/{1}} {
    \draw (\x,0.08) -- (\x,-0.08);
    \node[below,font=\scriptsize] at (\x,-0.10) {$\lbl$};
  }
  \node[font=\footnotesize] at (3,-0.65) {the sample space $\Omega = [0,1]$ (uniform length = probability)};
\end{tikzpicture}
\caption{The typewriter sequence. Each colored block shows the interval on which
the indicator $X_n$ equals $1$; elsewhere it is $0$. The sliding bump sweeps across
$[0,1]$ and shrinks on each pass. At any fixed outcome $\omega$, the bump sweeps past
infinitely many times, so $X_n(\omega)$ oscillates between $0$ and $1$ forever -- yet
the bump's width (the probability of error $\P(X_n \ne 0)$) drops to zero as $k \to \infty$.}
\label{ch22:fig-typewriter}
\end{figure}

\textbf{In probability it converges to $0$.} The variable $X_n$ equals $1$ only on a
block of width $1/k$ (where $k$ is the pass containing index $n$), so for any
$\eps\in(0,1)$,
\[
  \P(\abs{X_n - 0}\ge\eps) = \P(X_n = 1) = \tfrac{1}{k} \to 0
  \quad\text{as } n\to\infty,
\]
since the pass number $k$ grows without bound. Hence $X_n \pto 0$.

\textbf{Almost surely it converges nowhere.} Fix any point $\omega\in[0,1]$. In every
pass $k$, exactly one block contains $\omega$, so in every pass there is an index $n$
with $X_n(\omega) = 1$; and there are also indices with $X_n(\omega) = 0$. Thus the
numerical sequence $X_n(\omega)$ contains infinitely many $1$'s and infinitely many
$0$'s -- it oscillates forever and has no limit. The convergence fails for \emph{every}
$\omega$, the worst possible failure. So $X_n$ does \emph{not} converge to $0$ almost
surely (it fails on a set of probability $1$, not $0$).

\begin{happeningbox}
This single example proves $\pto$ does not imply $\asto$. The mechanism is exactly the
``moving spotlight'': at time $n$ the bad set (where $X_n=1$) is small, but it
\emph{relocates} each step, so no individual $\omega$ is spared in the long run. Almost
sure convergence forbids exactly this kind of perpetual relocation -- it demands the
failures eventually stop happening along each path. The Borel--Cantelli connection:
$\sum_n \P(X_n=1) = \sum_k k\cdot\frac1k = \sum_k 1 = \infty$, so the summable-tail
test is \emph{not} satisfied, consistent with a.s. failure.
\end{happeningbox}

\begin{walkbox}
\textbf{Pause and think: can a subsequence escape the typewriter trap?}
The full typewriter sequence $X_n$ oscillates forever at every single point $\omega$.
Now consider what happens if you throw away almost all terms and keep only a sparse
\emph{subsequence}.

Suppose you pick only the first block of each pass: $Y_k = \ind_{[0, 1/k]}$ for
$k = 1, 2, 3, \dots$. Does $Y_k$ converge almost surely to $0$?
Take any $\omega \in (0,1]$. As soon as $k > 1/\omega$, the interval $[0, 1/k]$ no
longer contains $\omega$, so $Y_k(\omega) = 0$ for all large $k$. The only point that
never settles is the single boundary point $\omega = 0$. But $\{\omega = 0\}$ has
probability zero! So $Y_k \asto 0$.

By thinning the sequence, the spotlight stopped sweeping back and forth. This is a
general mathematical theorem (due to Frigyes Riesz): \emph{every sequence that
converges in probability contains a subsequence that converges almost surely.}
\end{walkbox}

\section{Convergence in \texorpdfstring{$L^p$}{Lp} (in mean)}
\label{ch22:lp}

\begin{definition}[Convergence in $L^p$]
\label{ch22:def-lp}
Fix $p\ge1$. We say $X_n \to X$ \emph{in $L^p$} (or in $p$th mean), written
$X_n \xrightarrow{L^p} X$, if $\E\abs{X_n}^p<\infty$ for all $n$ and
\[
  \lim_{n\to\infty}\E\abs{X_n - X}^p = 0.
\]
The case $p=1$ is \emph{convergence in mean}; $p=2$ is \emph{convergence in mean
square} (or in quadratic mean), ubiquitous in statistics and signal processing.
\end{definition}

This mode measures error with an \emph{integrated} ruler rather than a probabilistic
one. It is sensitive to how \emph{large} the error can be, not merely how \emph{often}
it is nonnegligible. That sensitivity is exactly why it can fail when convergence in
probability holds: a rare error that is enormous can keep the average error large while
the probability of any error stays small.

\subsection{Convergence in $L^p$ implies convergence in probability}
\label{ch22:lp-implies-p}

\begin{theorem}
\label{ch22:thm-lp-p}
If $X_n \xrightarrow{L^p} X$ for some $p\ge1$, then $X_n \pto X$.
\end{theorem}

\begin{proof}
This is a one-line consequence of Markov's inequality applied to $\abs{X_n - X}^p$.
For any $\eps>0$,
\[
  \P(\abs{X_n - X}\ge\eps)
  = \P(\abs{X_n - X}^p \ge \eps^p)
  \le \frac{\E\abs{X_n - X}^p}{\eps^p}.
\]
The right side tends to $0$ by hypothesis, so the left side does too.
\end{proof}

\begin{connectionbox}
Notice the engine here is Markov's inequality, the most elementary concentration bound
(Chapter~19). A pattern worth absorbing: \emph{control of a moment
(an average) buys control of a probability (a tail).} The whole subject of
concentration of measure is built on upgrading this trade. The factor $1/\eps^p$ also
explains why larger $p$ gives a stronger statement -- a higher moment penalizes large
deviations more harshly, so $L^p$ convergence for big $p$ is harder to achieve and
implies more.
\end{connectionbox}

\subsection{In probability but not in $L^p$: the rare giant}
\label{ch22:giant}

The converse fails, and the reason is a single rare but huge excursion.

\begin{example}[A spike that is unlikely but enormous]
\label{ch22:ex-spike}
On $\Omega=[0,1]$ with uniform probability, define
\[
  X_n(\omega) =
  \begin{cases}
    n & \text{if } \omega \in (0, 1/n],\\
    0 & \text{otherwise.}
  \end{cases}
\]
Then $X_n \pto 0$: for any $\eps>0$ and $n$ large, $\P(\abs{X_n}\ge\eps)=\P(X_n=n)=
1/n \to 0$. But the mean does not vanish:
\[
  \E\abs{X_n - 0} = \E X_n = n\cdot\tfrac1n = 1 \quad\text{for all } n,
\]
so $X_n$ does \emph{not} converge to $0$ in $L^1$ (and a fortiori not in $L^2$, where
$\E X_n^2 = n^2\cdot\frac1n = n \to\infty$). The probability of error vanishes, but
the \emph{magnitude} when it occurs grows just fast enough to keep the average error
pinned at $1$. The area of the spike refuses to shrink.
\end{example}

\begin{figure}[H]
\centering
\begin{tikzpicture}
\begin{axis}[
  width=11.5cm, height=5.4cm,
  xlabel={$\omega$}, ylabel={$X_n(\omega)$},
  xmin=0, xmax=1, ymin=0, ymax=6.5,
  title style={font=\small},
  title={Rare giant spikes: height $n$, width $1/n$, area $\E[X_n]=1$},
  legend pos=north east, legend style={font=\small},
  grid=both, grid style={gray!18},
  every axis plot/.append style={very thick},
]
% n=2: height 2 on (0,1/2]
\addplot[sourceblue, const plot] coordinates {(0,2) (0.5,2) (0.5,0) (1,0)};
\addlegendentry{$X_2$ (height $2$, width $1/2$)}
% n=4: height 4 on (0,1/4]
\addplot[targetgreen, const plot] coordinates {(0,4) (0.25,4) (0.25,0) (1,0)};
\addlegendentry{$X_4$ (height $4$, width $1/4$)}
% n=6: height 6 on (0,1/6]
\addplot[densityred, const plot] coordinates {(0,6) (0.16667,6) (0.16667,0) (1,0)};
\addlegendentry{$X_6$ (height $6$, width $1/6$)}
\end{axis}
\end{tikzpicture}
\caption{The rare giant of Example~\ref{ch22:ex-spike}. Each spike is taller and
narrower than the last, but the rectangle's area (its height $n$ times its width $1/n$)
is always exactly $1$. So $\E X_n=1$ for every $n$ even as $\P(X_n\ne0)=1/n\to0$:
convergence in probability without convergence in $L^1$.}
\label{ch22:fig-spike}
\end{figure}

\begin{trapbox}
A natural-sounding but \textbf{false} belief: ``if $X_n\pto X$ then surely the means
converge, $\E X_n \to \E X$.'' Example~\ref{ch22:ex-spike} kills it: $X_n\pto0$ but
$\E X_n = 1 \not\to 0$. Convergence in probability says nothing about averages
unless you add a control on the tails. The standard such control is \emph{uniform
integrability} (a condition guaranteeing that probability mass cannot escape into
arbitrarily huge, rare outliers); under it, $\pto$ does upgrade to $L^1$ and the means
converge. Without some moment control, never pass a limit inside an expectation just
because things converge in probability.
\end{trapbox}

\section{Convergence in distribution}
\label{ch22:dist}

The final mode is the most permissive and, despite that, the most famous, because the
Central Limit Theorem is a statement of exactly this type.

\begin{definition}[Convergence in distribution]
\label{ch22:def-dist}
Let $X_n$ have cumulative distribution functions $F_n$ and let $X$ have CDF $F$. We
say $X_n$ \emph{converges in distribution} (or in law, or weakly) to $X$, written
$X_n \dto X$, if
\[
  \lim_{n\to\infty} F_n(x) = F(x)
\]
at every point $x$ where $F$ is continuous.
\end{definition}

Two features of this definition repay close reading. First, it speaks only of the
\emph{distributions} $F_n, F$ -- the random variables themselves may live on entirely
different probability spaces, may have nothing to do with one another as functions.
Second, the qualifier ``at every continuity point of $F$'' is not a technicality you
may ignore; it is essential.

\begin{whybox}
Why exempt the discontinuity points of the limit? Consider $X_n$ deterministically
equal to $1/n$, converging in every sane sense to the constant $0$. The limiting CDF
$F$ of the constant $0$ jumps from $0$ to $1$ at $x=0$, so $F(0)=1$. But
$F_n(0)=\P(1/n\le0)=0$ for every $n$, so $F_n(0)=0\not\to1=F(0)$ \emph{at the jump}.
Everywhere else $F_n\to F$ perfectly. If we demanded convergence at the jump too, we
would absurdly conclude $1/n\not\to0$. Exempting continuity points repairs this. The
moral: convergence in distribution is about the \emph{shape} of the distribution
settling down, and a CDF's value exactly at a jump is an arbitrary bookkeeping choice
that the shape should not be hostage to.
\end{whybox}

\begin{figure}[H]
\centering
\begin{tikzpicture}
\begin{axis}[
  width=11.5cm, height=6.2cm,
  xlabel={$x$}, ylabel={CDF},
  domain=-3.5:3.5, samples=160,
  ymin=-0.05, ymax=1.08, xmin=-3.5, xmax=3.5,
  title style={font=\small},
  title={Convergence in distribution: cumulative distributions $F_n \to F$},
  legend pos=north west, legend cell align=left,
  legend style={font=\footnotesize},
  every axis plot/.append style={very thick},
  grid=both, grid style={gray!18},
]
% Logistic curves used as smooth stand-ins for normal CDFs (pgfplots has no erf).
\addplot[densityred, dashed] {1/(1+exp(-1.702*x))};
\addlegendentry{limit CDF $F$ (standard normal)}
\addplot[sourceblue] {1/(1+exp(-1.0*(x-0.4)))};
\addlegendentry{$F_n$, small $n$ (off in shape and location)}
\addplot[targetgreen] {1/(1+exp(-1.45*(x-0.1)))};
\addlegendentry{$F_n$, larger $n$ (closing in)}
\end{axis}
\end{tikzpicture}
\caption{Convergence in distribution is convergence of CDFs (at continuity points).
As $n$ grows, the curve $F_n$ relaxes onto the limiting curve $F$. Nothing is claimed
about the underlying random values lining up -- only the cumulative shapes do.
(Logistic curves are used here as smooth stand-ins for normal CDFs.)}
\label{ch22:fig-cdf-conv}
\end{figure}

\subsection{Equivalent formulations: the Portmanteau theorem}
\label{ch22:portmanteau}

Convergence in distribution has several equivalent characterizations; knowing them is what
makes it usable in proofs and calculations. The collection is traditionally called the \emph{Portmanteau theorem}
(from the French for a collection of diverse items gathered together).

\begin{theorem}[Portmanteau, selected equivalences]
\label{ch22:thm-portmanteau}
The following are equivalent to $X_n \dto X$:
\begin{enumerate}
\item $F_n(x)\to F(x)$ at every continuity point of $F$.
\item $\E\, g(X_n) \to \E\, g(X)$ for every bounded continuous function $g$.
\item $\E\, g(X_n) \to \E\, g(X)$ for every bounded Lipschitz function $g$ (functions
      whose rate of change is bounded by a constant, $\abs{g(x)-g(y)}\le K\abs{x-y}$).
\item For the characteristic functions $\varphi_{X_n}(t) = \E[e^{itX_n}]$ (the Fourier
      transforms that uniquely determine each distribution), $\varphi_{X_n}(t)\to\varphi_X(t)$
      for every $t$ (this is L\'evy's continuity theorem, guaranteeing that pointwise
      convergence of characteristic functions captures convergence of distributions).
\end{enumerate}
\end{theorem}

Form (2) is the conceptual core: convergence in distribution means \emph{every
bounded continuous measurement} you could take has its expectation converge. Form (4) provides
the standard analytical tool for proofs: it is how the Central Limit Theorem is proved,
by showing the characteristic function of the standardized sum converges to
$e^{-t^2/2}$, the characteristic function of the standard normal (see Chapter~21).

\begin{connectionbox}
The phrase ``test against bounded continuous functions'' is the same idea that defines
weak convergence of measures throughout analysis, and the same spirit as ``test
against smooth functions'' in the theory of distributions (generalized functions). If
you have met the companion calculus book's treatment of integration as pairing a
function against test functions, form (2) will feel like home: a distribution is
known by how it integrates against probes.
\end{connectionbox}

\section{The implication diagram and intuitive filters}
\label{ch22:diagram}

We now assemble the full hierarchy of stochastic convergence. The four modes
act as an intuitive cascade of filters, each testing a different property of the
sequence of random variables:
\begin{itemize}
  \item \textbf{Pathwise lock-in filter (Almost Sure):} Evaluates each sample path $\omega \mapsto X_n(\omega)$ individually. Tests whether the path enters an $\eps$-tube around $X(\omega)$ and stays trapped inside for all large $n$.
  \item \textbf{Energy / outlier penalty filter ($L^p$ Mean):} Evaluates the overall magnitude of deviations, $\E\abs{X_n - X}^p$. Harshly penalizes rare but extreme spikes.
  \item \textbf{Flickering mass filter (In Probability):} Takes a snapshot of the whole sample space $\Omega$ at time $n$. Tests whether the total probability mass of points outside the $\eps$-tube vanishes as $n \to \infty$, even if individual points flicker in and out.
  \item \textbf{Distributional shape / CDF filter (In Distribution):} Discards the sample space entirely. Tests only whether the cumulative distribution functions $F_n(x)$ align with $F(x)$ at continuity points.
\end{itemize}

Figure~\ref{ch22:fig-diagram} lays out the resulting hierarchy, the forward implications, and the counterexamples demonstrating where reverse arrows fail.

\begin{figure}[H]
\centering
\begin{tikzpicture}[
  scale=0.88,
  box/.style={draw=#1, line width=1pt, rounded corners=4pt, align=center,
              font=\small, inner sep=5pt},
  imp/.style={-{Stealth[length=3mm]}, very thick},
  cex/.style={draw=trapred!70, fill=trapred!8, rounded corners=3pt, font=\scriptsize, align=center, inner sep=3.5pt},
  bridge/.style={draw=targetgreen!70, fill=targetgreen!8, rounded corners=3pt, font=\scriptsize, align=center, inner sep=3.5pt}
]
  % Top lane: Strongest modes
  \node[box=structurepurple, fill=structurepurple!12, minimum width=4.3cm] (as) at (0.2,4.65) {
    \textbf{Almost Surely ($X_n \asto X$)}\\[2pt]
    \footnotesize \textit{Pathwise lock-in filter}\\
    \scriptsize $\P(\lim_{n\to\infty} X_n(\omega) = X(\omega)) = 1$
  };

  \node[box=quotientorange, fill=quotientorange!12, minimum width=4.3cm] (lp) at (5.4,4.65) {
    \textbf{In $L^p$ Mean ($X_n \xrightarrow{L^p} X$)}\\[2pt]
    \footnotesize \textit{Energy / outlier filter}\\
    \scriptsize $\lim_{n\to\infty} \E\abs{X_n - X}^p = 0$
  };

  % Middle lane: In Probability
  \node[box=sourceblue, fill=sourceblue!12, minimum width=4.8cm] (prob) at (2.8,2.0) {
    \textbf{In Probability ($X_n \pto X$)}\\[2pt]
    \footnotesize \textit{Flickering mass filter}\\
    \scriptsize $\lim_{n\to\infty} \P(\abs{X_n - X} \ge \eps) = 0$
  };

  % Bottom lane: In Distribution
  \node[box=targetgreen, fill=targetgreen!12, minimum width=4.8cm] (dist) at (2.8,-0.2) {
    \textbf{In Distribution ($X_n \dto X$)}\\[2pt]
    \footnotesize \textit{Distributional shape / CDF filter}\\
    \scriptsize $\lim_{n\to\infty} F_n(x) = F(x)$ at continuity pts
  };

  % Forward implication arrows
  \draw[imp, structurepurple] (as.south) -- (prob.125)
    node[midway, left=2pt, font=\scriptsize, structurepurple] {Tail bounds};
  \draw[imp, quotientorange] (lp.south) -- (prob.55)
    node[midway, right=2pt, font=\scriptsize, quotientorange] {Markov inequality};
  \draw[imp, sourceblue] (prob.south) -- (dist.north)
    node[midway, right=3pt, font=\scriptsize, sourceblue] {Portmanteau};

  % Counterexample callouts
  \node[cex] (typewriter) at (-2.6,2.0) {
    \textbf{No reverse: Typewriter}\\
    sliding bump hits all $\omega$ i.o.;\\
    $\pto 0$ holds, $\asto 0$ fails
  };
  \draw[trapred, thick, dashed, -{Stealth[length=2.5mm]}] (typewriter.east) -- (prob.west);

  \node[cex] (spike) at (8.5,2.0) {
    \textbf{No reverse: Giant spike}\\
    $X_n = n^2 \ind_{[0,1/n]}$;\\
    $\pto 0$ holds, $\E[X_n]=n \to \infty$
  };
  \draw[trapred, thick, dashed, -{Stealth[length=2.5mm]}] (spike.west) -- (prob.east);

  \node[cex] (coin) at (9.1,-0.2) {
    \textbf{No reverse: Opposite coin}\\
    $X_n = 1-X$ with $X \sim \Bern(\tfrac12)$;\\
    $\dto$ holds, but $\P(\abs{X_n-X}=1)=1$
  };
  \draw[trapred, thick, dashed, -{Stealth[length=2.5mm]}] (coin.west) -- (dist.east);

  \node[bridge] (const) at (-2.6,-0.2) {
    \textbf{Bridge for constants:}\\
    If $X = c$ is constant,\\
    $X_n \dto c \;\Longleftrightarrow\; X_n \pto c$
  };
  \draw[targetgreen!80!black, thick, Stealth-Stealth] (const.east) -- (dist.west);

\end{tikzpicture}
\caption{The hierarchy of stochastic convergence. Almost sure and $L^p$ each imply
convergence in probability, which implies convergence in distribution. None of the
reverse implications hold in general: the typewriter sequence shows $\pto \not\Rightarrow \asto$;
rare giant spikes show $\pto \not\Rightarrow L^p$; and opposite coin flips show
$\dto \not\Rightarrow \pto$. For a constant limit $c$, however, $\dto c$ upgrades to $\pto c$.}
\label{ch22:fig-diagram}
\end{figure}

Let us nail down each arrow and each \emph{non}-arrow.

\subsection{Almost sure \texorpdfstring{$\Rightarrow$}{=>} in probability}
\label{ch22:as-imp-p}

\begin{theorem}
\label{ch22:thm-as-p}
If $X_n\asto X$ then $X_n\pto X$.
\end{theorem}

\begin{proof}
Fix $\eps>0$ and define the tail event $A_N = \bigcup_{n\ge N}\set{\abs{X_n-X}\ge\eps}$,
the event that some deviation of size $\eps$ occurs at time $N$ or later. These events
decrease: $A_1 \supseteq A_2 \supseteq \cdots$. Their intersection is the event
``$\abs{X_n - X}\ge\eps$ infinitely often,'' which has probability $0$ because almost
sure convergence forbids infinitely-often large deviations
(Section~\ref{ch22:as-tube}). By continuity of probability from above,
$\P(A_N)\to\P(\bigcap_N A_N)=0$. But $\set{\abs{X_N-X}\ge\eps}\subseteq A_N$, so
$\P(\abs{X_N-X}\ge\eps)\le\P(A_N)\to0$. Hence $X_n\pto X$.
\end{proof}

The typewriter sequence (Section~\ref{ch22:typewriter}) shows the reverse fails: it
converges in probability but not almost surely.

\subsection{In probability \texorpdfstring{$\Rightarrow$}{=>} in distribution}
\label{ch22:p-imp-d}

\begin{theorem}
\label{ch22:thm-p-d}
If $X_n\pto X$ then $X_n\dto X$.
\end{theorem}

\begin{proof}[Proof sketch]
Let $x$ be a continuity point of $F$. For any tolerance $\delta>0$, we split on whether
$X_n$ is within $\delta$ of $X$:
\[
  F_n(x)=\P(X_n\le x)
  \le \P(X\le x+\delta) + \P(\abs{X_n - X} > \delta).
\]
Why does this set inclusion hold? Whenever $X_n \le x$, either $X \le x+\delta$ or
$X > x+\delta$. But if $X > x+\delta$ and $X_n \le x$, the gap between them satisfies
$X - X_n > (x+\delta) - x = \delta$, which forces $\abs{X_n - X} > \delta$.

Now send $n\to\infty$. The probability $\P(\abs{X_n - X} > \delta)$ vanishes by
hypothesis ($X_n \pto X$), leaving
\[
  \limsup_{n\to\infty} F_n(x) \le F(x+\delta).
\]
Running the exact same argument symmetrically with the roles reversed gives
$\liminf_{n\to\infty} F_n(x) \ge F(x-\delta)$. Finally, let $\delta \downarrow 0$:
because $x$ is a continuity point of $F$, both $F(x+\delta)$ and $F(x-\delta)$ converge
to $F(x)$, squeezing $F_n(x) \to F(x)$.
\end{proof}

\begin{trapbox}
The reverse, ``in distribution $\Rightarrow$ in probability,'' is false in a way that
beginners trip on constantly. Let $Z\sim\Normal(0,1)$ and set $X_n = (-1)^n Z$. By
symmetry of the normal, every $X_n$ has the \emph{same} distribution as $Z$, so
immediately $X_n\dto Z$. But $\abs{X_n - Z}$ is either $0$ (even $n$) or $2\abs{Z}$
(odd $n$), so for odd $n$, $\P(\abs{X_n - Z}\ge\eps)=\P(2\abs Z\ge\eps)\not\to0$. The
distributions agree perfectly while the \emph{values} stay far apart. Distribution
convergence sees the shape, never the alignment.
\end{trapbox}

\subsection{One exception: convergence in distribution to a constant}
\label{ch22:to-constant}

There is exactly one important case where the weak mode upgrades for free.

\begin{theorem}
\label{ch22:thm-const}
If $c$ is a constant, then $X_n\dto c$ if and only if $X_n\pto c$.
\end{theorem}

\begin{proof}
($\Leftarrow$) is Theorem~\ref{ch22:thm-p-d}. For ($\Rightarrow$): the limit variable
$c$ is deterministic, so its cumulative distribution is a pure step: $F(t) = 0$ for
$t < c$ and $F(t) = 1$ for $t \ge c$. For any $\eps>0$,
\[
  \P(\abs{X_n - c}\ge\eps)
  = \P(X_n\le c-\eps) + \P(X_n\ge c+\eps).
\]
Because $\{X_n \ge c+\eps\} \subseteq \{X_n > c+\eps/2\}$, the second term is bounded by
$1 - F_n(c+\eps/2)$. Hence
\[
  \P(\abs{X_n - c}\ge\eps) \le F_n(c-\eps) + \big(1 - F_n(c+\tfrac{\eps}{2})\big).
\]
Why did we shift to $c+\eps/2$? Because the limit CDF jumps at $c$ alone; any point
strictly to the right, like $c+\eps/2$, is a continuity point of $F$ where $F(c+\eps/2) = 1$.
Similarly, $c-\eps < c$ is a continuity point with $F(c-\eps) = 0$. Since $X_n \dto c$,
$F_n(c-\eps) \to 0$ and $F_n(c+\eps/2) \to 1$, driving the right-hand sum to zero.
\end{proof}

This little theorem is the secret ingredient in Slutsky's theorem below, and it is
why the Weak Law of Large Numbers (convergence in probability to the constant $\mu$)
can equivalently be phrased as convergence in distribution to $\mu$.

\begin{recoverybox}
Pause and consolidate the four facts you now own:
\begin{itemize}
\item $\asto \Rightarrow \pto$, but not conversely (typewriter).
\item $\xrightarrow{L^p} \Rightarrow \pto$, but not conversely (rare giant spike).
\item $\pto \Rightarrow \dto$, but not conversely (the $(-1)^nZ$ trap).
\item a.s. and $L^p$ imply neither each other; and to a \emph{constant} limit, $\dto$
      and $\pto$ coincide.
\end{itemize}
Everything else in the chapter is built from these.
\end{recoverybox}

\subsection{Almost sure and \texorpdfstring{$L^p$}{Lp} do not imply each other}
\label{ch22:as-vs-lp}

These two ``strong'' modes are siblings, not parent and child; neither implies the
other.

\begin{itemize}
\item \emph{$L^p$ but not a.s.}: the typewriter sequence itself. With the sliding
      block of width $1/k$ from Section~\ref{ch22:typewriter}, we have
      $\E\abs{X_n}^p = 1/k \to 0$, so $X_n\to0$ in $L^p$, yet (as shown) $X_n$
      converges almost surely nowhere.
\item \emph{a.s.\ but not $L^p$}: the rare giant of Example~\ref{ch22:ex-spike}
      (Figure~\ref{ch22:fig-spike}). For a
      fixed $\omega\in(0,1]$ we have $\omega>1/n$ for all large $n$, hence
      $X_n(\omega)=0$ eventually; so $X_n\asto0$. Yet $\E X_n = 1\not\to0$, so it
      fails in $L^1$.
\end{itemize}

Both ``strong'' modes imply convergence in probability, but they control \emph{different
things} -- one the long-run behavior of each path, the other the average magnitude of
the error -- and a sequence can ace one while flunking the other.

\section{The continuous mapping theorem}
\label{ch22:cmt}

Convergence would be of limited use if we could not push it through functions. The continuous mapping theorem says: a continuous transformation respects three of the modes: almost sure, in probability, and in distribution.

\begin{theorem}[Continuous mapping]
\label{ch22:thm-cmt}
Let $g$ be a function continuous on a set $C$ with $\P(X\in C)=1$ (continuous at every
value $X$ actually takes). Then:
\[
  X_n\asto X \;\Rightarrow\; g(X_n)\asto g(X),\qquad
  X_n\pto X \;\Rightarrow\; g(X_n)\pto g(X),\qquad
  X_n\dto X \;\Rightarrow\; g(X_n)\dto g(X).
\]
\end{theorem}

\begin{whybox}
Why continuity is exactly the right hypothesis. Continuity of $g$ means: inputs that
are close get mapped to outputs that are close. Each mode of convergence is
a statement about inputs getting close (pointwise, in probability, in distribution).
Continuity is precisely the property that transfers ``inputs close'' to ``outputs close,'' so it is no surprise it transfers these three modes. For the distributional case,
the cleanest proof uses Portmanteau form (2): for bounded continuous $h$, the
composition $h\circ g$ is bounded continuous, so $\E\,h(g(X_n))=\E\,(h\circ g)(X_n)\to
\E\,(h\circ g)(X)=\E\,h(g(X))$. The result falls right out of testing against probes.
\end{whybox}

\begin{example}[From the CLT to a chi-square limit]
\label{ch22:ex-chisq}
Suppose $\sqrt n(\bar X_n - \mu)/\sigma \dto Z\sim\Normal(0,1)$ (the CLT). The map
$g(x)=x^2$ is continuous, so by continuous mapping
\[
  \frac{n(\bar X_n - \mu)^2}{\sigma^2} \dto Z^2 \sim \chi^2_1.
\]
This is the seed of every chi-square test in statistics: square a normal limit and you
get a chi-square limit, automatically, with no new probability calculation.
\end{example}

\section{Slutsky's theorem}
\label{ch22:slutsky}

The continuous mapping theorem is about one converging sequence. Real statistics
constantly combines \emph{two}: an estimator that converges in distribution, and a
nuisance quantity (an estimated variance, say) that converges in probability to a
constant. Slutsky's theorem is the rule for combining them, and it may be the
most-used theorem in all of asymptotic statistics.

\begin{theorem}[Slutsky]
\label{ch22:thm-slutsky}
Suppose $X_n\dto X$ and $Y_n\pto c$ for a \emph{constant} $c$. Then:
\[
  X_n + Y_n \dto X + c,\qquad
  Y_n X_n \dto c\,X,\qquad
  \frac{X_n}{Y_n}\dto\frac{X}{c}\ \ (c\ne0).
\]
\end{theorem}

\begin{focusbox}
Two warnings that are the whole reason Slutsky is a named theorem and not a triviality.
\begin{enumerate}
\item The second sequence must converge to a \textbf{constant}, not merely in
      distribution. If $Y_n\dto Y$ for a nondegenerate $Y$, the conclusions are
      \textbf{false} in general -- because convergence in distribution carries no
      information about the \emph{joint} behavior of $X_n$ and $Y_n$, and a sum or
      product depends on the joint law. Constancy is what removes that dependence:
      a constant has no randomness to be jointly entangled with.
\item It is genuinely surprising that you may add a $\dto$ sequence and a $\pto$
      sequence at all, since $\dto$ alone does not respect addition. The constant limit
      is exactly the loophole that makes it legal.
\end{enumerate}
\end{focusbox}

\begin{proof}[Proof idea]
The trick is to convert $Y_n\pto c$ into something the joint limit can absorb. Write
$X_n + Y_n = (X_n + c) + (Y_n - c)$. The first piece converges in distribution to
$X+c$ (continuous mapping, shift by a constant). The second piece, $Y_n - c$,
converges in probability to $0$. A lemma says: if $A_n\dto A$ and $B_n\pto0$ then
$A_n + B_n\dto A$ (a vanishing-in-probability perturbation cannot change a
distributional limit -- intuitively, with high probability $B_n$ is negligible, and
on the rare event it is not, it contributes vanishing probability mass). Apply the
lemma. The product and quotient versions follow by the same perturbation lemma plus
continuous mapping.
\end{proof}

\begin{example}[The $t$-statistic: why you may plug in $s$ for $\sigma$]
\label{ch22:ex-tstat}
The cleanest application. By the CLT,
$\sqrt n(\bar X_n - \mu)/\sigma \dto \Normal(0,1)$. But in practice $\sigma$ is
unknown; we use the sample standard deviation $s_n$, which (by the LLN plus continuous
mapping, since $s_n^2\pto\sigma^2$ and the square root is continuous) satisfies
$s_n\pto\sigma$, hence $\sigma/s_n\pto1$. Slutsky's product rule then gives
\[
  \frac{\sqrt n(\bar X_n - \mu)}{s_n}
  = \underbrace{\frac{\sqrt n(\bar X_n - \mu)}{\sigma}}_{\dto\,\Normal(0,1)}
    \cdot
    \underbrace{\frac{\sigma}{s_n}}_{\pto\,1}
  \;\dto\; \Normal(0,1).
\]
So replacing the unknown $\sigma$ by its estimate $s_n$ does not change the limiting
distribution. This is precisely why large-sample $t$-tests and Wald confidence
intervals are valid even though we never know the true variance.
\end{example}

\begin{applicationbox}
\textbf{A/B testing and Wald intervals.} When you compute a $95\%$ confidence interval
$\hat\theta \pm 1.96\,\widehat{\mathrm{SE}}$ for a conversion-rate difference, two
asymptotic facts are quietly cooperating: the estimator $\hat\theta$ is asymptotically
normal (CLT, a $\dto$ statement) and the estimated standard error
$\widehat{\mathrm{SE}}$ is consistent for the true SE (LLN, a $\pto$ statement to a
constant). Slutsky fuses them so that the studentized statistic
$(\hat\theta-\theta)/\widehat{\mathrm{SE}}$ is asymptotically standard normal, which
is what makes the $\pm1.96$ legitimate. Every dashboard that reports a $p$-value from
a $z$- or $t$-statistic relies on Slutsky's theorem.
\end{applicationbox}

\section{The delta method}
\label{ch22:delta}

The CLT tells you the asymptotic distribution of an average. But you usually care
about a \emph{nonlinear function} of an average -- a log-odds, a ratio, a variance, or a
correlation. The delta method is the rule that transports asymptotic normality through
a smooth function. It combines a first-order Taylor expansion with Slutsky's theorem.

\begin{theorem}[Delta method]
\label{ch22:thm-delta}
Suppose $\sqrt n\,(T_n - \theta) \dto \Normal(0,\sigma^2)$ and $g$ is differentiable
at $\theta$ with $g'(\theta)\ne0$. Then
\[
  \sqrt n\,\big(g(T_n) - g(\theta)\big) \dto \Normal\!\big(0,\ \sigma^2\,[g'(\theta)]^2\big).
\]
\end{theorem}

\begin{intuitionbox}
Zoom in on the graph of $g$ near $\theta$. Because $T_n$ is concentrating on $\theta$
(its fluctuations are of order $1/\sqrt n$), only the behavior of $g$ in a tiny
neighborhood of $\theta$ matters. On that tiny scale $g$ looks like its tangent line,
slope $g'(\theta)$. A straight line of slope $m$ turns a $\Normal(0,\sigma^2)$ wiggle
into a $\Normal(0,m^2\sigma^2)$ wiggle (variance scales by the square of the slope).
That is the entire content: \emph{locally linearize, then scale the variance by the
slope squared.}
\end{intuitionbox}

\begin{figure}[H]
\centering
\begin{tikzpicture}
\begin{axis}[
  width=11cm, height=5.6cm,
  xlabel={$t$}, ylabel={$g(t)$},
  domain=0.2:3, samples=120,
  xmin=0.2, xmax=3, ymin=-1.4, ymax=1.4,
  title style={font=\small},
  title={The delta method: local linearization $g(t) \approx g(\theta) + g'(\theta)(t-\theta)$},
  legend style={font=\footnotesize, at={(0.98,0.06)}, anchor=south east, column sep=4pt},
  legend cell align=left,
  grid=both, grid style={gray!18},
  every axis plot/.append style={very thick},
]
\addplot[sourceblue] {ln(x)};
\addlegendentry{$g(t)=\ln t$}
\addplot[densityred, dashed, domain=0.6:1.8] {ln(1.2) + (1/1.2)*(x-1.2)};
\addlegendentry{tangent line (slope $g'(\theta)$)}
\addplot[gray!60, dotted, thick] coordinates {(1.2, -1.4) (1.2, 0.18232)};
\addplot[only marks, mark=*, mark size=2pt, black] coordinates {(1.2, 0.18232)};
\node[above left, font=\footnotesize] at (axis cs:1.2, 0.18232) {$(\theta, g(\theta))$};
\node[above right, font=\footnotesize] at (axis cs:1.2, -1.35) {$\theta = 1.2$};
\end{axis}
\end{tikzpicture}
\caption{The delta method as local linearization. Where $T_n$ concentrates (near
$\theta$), the smooth function $g$ is indistinguishable from its tangent line. The
tangent's slope $g'(\theta)$ is the only feature of $g$ that survives into the limiting
variance, which scales by $[g'(\theta)]^2$.}
\label{ch22:fig-delta}
\end{figure}

\begin{proof}
A first-order Taylor expansion of $g$ about $\theta$ gives, for some $\theta_n^\star$
between $T_n$ and $\theta$,
\[
  g(T_n) = g(\theta) + g'(\theta_n^\star)\,(T_n - \theta).
\]
Multiply by $\sqrt n$:
\[
  \sqrt n\,\big(g(T_n) - g(\theta)\big) = g'(\theta_n^\star)\cdot \sqrt n\,(T_n-\theta).
\]
Now use the machinery of this chapter. Since $\sqrt n(T_n-\theta)\dto\Normal(0,\sigma^2)$,
the unscaled fluctuation $T_n-\theta$ is of order $1/\sqrt n$ and so $T_n\pto\theta$.
The intermediate point $\theta_n^\star$ is trapped between $T_n$ and $\theta$, so
$\theta_n^\star\pto\theta$, and by continuity of $g'$ at $\theta$,
$g'(\theta_n^\star)\pto g'(\theta)$, a constant. We now have a product of a
$\pto$-constant factor and a $\dto$-normal factor -- exactly Slutsky's product rule:
\[
  \sqrt n\,\big(g(T_n)-g(\theta)\big)
  \;\dto\; g'(\theta)\cdot\Normal(0,\sigma^2)
  = \Normal\!\big(0,[g'(\theta)]^2\sigma^2\big).\qedhere
\]
\end{proof}

\begin{example}[Variance-stabilizing a proportion]
\label{ch22:ex-arcsine}
Let $\hat p_n$ be the sample proportion of successes in $n$ Bernoulli$(p)$ trials. The
CLT gives $\sqrt n(\hat p_n - p)\dto\Normal(0,\,p(1-p))$ -- note the limiting variance
$p(1-p)$ depends on the unknown $p$, an annoyance. Apply the delta method with
$g(x)=\arcsin(\sqrt x)$, whose derivative is $g'(x)=\frac{1}{2\sqrt{x(1-x)}}$. Then
$[g'(p)]^2\,p(1-p)=\frac{1}{4p(1-p)}\cdot p(1-p)=\frac14$, a constant free of $p$:
\[
  \sqrt n\big(\arcsin\sqrt{\hat p_n}-\arcsin\sqrt p\big)\dto\Normal\!\big(0,\tfrac14\big).
\]
The arcsine-root transform \emph{stabilizes the variance}. This is exactly why
classical statistics uses it before plotting or modeling proportions: the delta method
told us which nonlinear lens makes the noise uniform.
\end{example}

\begin{walkbox}
\textbf{Pause and think: what happens when the derivative vanishes?}
Theorem~\ref{ch22:thm-delta} explicitly requires $g'(\theta) \ne 0$. What if you want the
asymptotic distribution of $g(T_n) = (T_n - \theta)^2$, where $g'(\theta) = 0$?

Because the first-order Taylor term vanishes, local linearity no longer governs the
fluctuations. We must look to the quadratic term:
$g(T_n) - g(\theta) \approx \frac{1}{2} g''(\theta)\,(T_n - \theta)^2$.
If $\sqrt n\,(T_n - \theta) \dto Z \sim \Normal(0,\sigma^2)$, then squaring both sides
gives
\[
  n\,\big(g(T_n) - g(\theta)\big) \approx \tfrac{1}{2} g''(\theta)\,\big(\sqrt n(T_n-\theta)\big)^2
  \;\dto\; \tfrac{1}{2} g''(\theta)\,\sigma^2\,Z^2
  \sim \tfrac{1}{2} g''(\theta)\,\sigma^2\,\chi^2_1.
\]
Notice two things changed: the scaling rate accelerated from $\sqrt n$ to $n$, and the
limit transformed from a normal bell curve into a chi-square distribution! The delta
method didn't break; calculus simply advanced to the next non-vanishing Taylor term.
\end{walkbox}

\begin{connectionbox}
The single-variable delta method is a Taylor expansion; the multivariable version is a
Taylor expansion with a \emph{gradient}. If $\sqrt n(T_n-\theta)\dto\Normal(0,\Sigma)$
for a vector parameter, then $\sqrt n(g(T_n)-g(\theta))\dto\Normal(0,\nabla g^\top
\Sigma\,\nabla g)$. The scalar slope-squared $[g'(\theta)]^2$ becomes the quadratic
form $\nabla g^\top\Sigma\nabla g$ -- precisely the linear-algebra object the companion
linear algebra book studies when a covariance matrix is sandwiched between a vector and
its transpose. Local linearization is the bridge: calculus gives the gradient, linear
algebra propagates the covariance.
\end{connectionbox}

\section{Tying in the Law of Large Numbers and the CLT}
\label{ch22:lln-clt}

The major asymptotic theorems of the previous two chapters correspond directly to these modes.

Let $X_1,X_2,\dots$ be i.i.d.\ with mean $\mu$ and finite variance $\sigma^2$, and
$\bar X_n = \frac1n\sum_{i=1}^n X_i$.

\textbf{The Weak Law of Large Numbers} is convergence in probability:
\[
  \bar X_n \pto \mu.
\]
The sample mean is, at each large $n$, very likely close to the truth. Snapshot the
crowd; almost all of it is near $\mu$.

\textbf{The Strong Law of Large Numbers} is almost sure convergence:
\[
  \bar X_n \asto \mu.
\]
Pick one infinite data stream; that particular running average lands on $\mu$ and
stays. This is genuinely stronger (as our hierarchy promised), which is why it takes
more work to prove and why the words ``weak'' and ``strong'' are well chosen.

\textbf{The Central Limit Theorem} is convergence in distribution:
\[
  \frac{\sqrt n\,(\bar X_n - \mu)}{\sigma} \dto \Normal(0,1).
\]
The CLT cannot be an a.s.\ or in-probability statement, and seeing why is illuminating.

\begin{happeningbox}
\textbf{Why the CLT must be distributional and nothing stronger.} The standardized
mean $\sqrt n(\bar X_n-\mu)/\sigma$ does \emph{not} converge as a sequence of numbers
at all -- it keeps wiggling, with standard deviation pinned at $1$ forever. There is
no number it homes in on, so almost-sure or in-probability convergence is simply not on
the table; the quantity never settles. What \emph{does} settle is its \emph{histogram}:
the shape of its distribution relaxes onto the bell curve. That is exactly what
convergence in distribution captures and the stronger modes cannot. The LLN squeezes
$\bar X_n$ onto a point; the CLT rescales the leftover wiggle by $\sqrt n$ and asks
about its \emph{shape}, which is a question only the weakest mode can answer. The two
theorems are looking at the same average through different magnifications.
\end{happeningbox}

\begin{figure}[H]
\centering
\begin{tikzpicture}
\begin{axis}[
  width=12cm, height=6cm,
  xlabel={$n$ (number of samples)},
  ylabel={running average $\bar X_n$},
  xmin=0, xmax=200, ymin=0.1, ymax=0.9,
  title style={font=\small},
  title={Strong Law of Large Numbers: pathwise convergence $\bar X_n \to \mu$},
  legend pos=north east, grid=both, grid style={gray!18},
  every axis plot/.append style={thick},
]
% true mean line
\addplot[densityred, dashed, domain=0:200, samples=2] {0.5};
\addlegendentry{$\mu = 0.5$}
% deterministic surrogates for three converging running averages
\addplot[sourceblue, samples=200, domain=1:200]
  {0.5 + 0.42*sin(deg(x*0.7))*exp(-x/55)/(1+0.03*x)};
\addlegendentry{path A}
\addplot[targetgreen, samples=200, domain=1:200]
  {0.5 - 0.38*cos(deg(x*0.9+1))*exp(-x/70)/(1+0.025*x)};
\addlegendentry{path B}
\addplot[quotientorange, samples=200, domain=1:200]
  {0.5 + 0.5*sin(deg(x*1.3+2))*exp(-x/45)/(1+0.04*x)};
\addlegendentry{path C}
\end{axis}
\end{tikzpicture}
\caption{The Law of Large Numbers visualized as almost sure convergence: three
individual sample paths of the running average, each homing in on and staying near
$\mu=0.5$. (Smooth deterministic curves stand in for true random paths.) The CLT lives
in the \emph{rate} at which the gap to $\mu$ closes -- it shrinks like $1/\sqrt n$, and
the rescaled gap is asymptotically normal.}
\label{ch22:fig-lln-paths}
\end{figure}

\begin{historybox}
The names ``weak'' and ``strong'' law, and the entire architecture of convergence modes,
crystallized over two centuries of struggle to make probability mathematically rigorous.
Jacob Bernoulli spent twenty years trying to prove that empirical frequencies in games
of chance must converge to true underlying probabilities (published posthumously in his
1713 \emph{Ars Conjectandi}). But Bernoulli had only finite combinatorics and binomial
bounds: he could only estimate tail probabilities at a fixed sample size $n$. That
naturally yielded convergence in probability (the Weak Law). He had no mathematical
language to speak of an infinite sequence of coin flips as a single completed object.

By 1909, \'Emile Borel possessed Henri Lebesgue's newly invented measure theory on the
unit interval $[0,1]$. Borel realized that an infinite sequence of coin flips could be
viewed as the binary expansion of a single real number $\omega \in [0,1]$ -- mapping an
entire trajectory through time to a single geometric point. This breakthrough allowed
him to prove the first Strong Law of Large Numbers (his ``normal numbers'' theorem,
showing almost every real number has balanced binary digits).

Yet during the 1920s, bitter controversy erupted over the very foundations of the
field. Richard von Mises insisted that probability could only be defined through
empirical frequencies in idealized physical experiments (his \emph{Kollektivs}),
rejecting abstract measure theory. Other mathematicians countered that von Mises's
rules were logically contradictory, while traditional analysts dismissed probability
as mere heuristic physics rather than pure mathematics.

Andrey Kolmogorov resolved the crisis in his 1933 masterpiece \emph{Grundbegriffe der
Wahrscheinlichkeitsrechnung}. By defining random variables as measurable functions on an
abstract probability space $(\Omega, \mathcal{F}, \P)$, Kolmogorov showed that probability
was simply measure theory with total measure $1$. Instantly, the four modes of
convergence snapped into sharp focus as the standard ways functions converge in
analysis: pointwise almost everywhere (almost sure), in measure (in probability), in
$L^p$ norm (in mean), and through integration against continuous test probes (in
distribution). Working in parallel, Evgeny Slutsky (1925) and Harald Cram\'er used this
clarity to place asymptotic statistics on rigorous foundations.

Today, while scalar convergence is thoroughly settled, the frontier has moved in two
directions: to infinite-dimensional spaces (empirical process theory and Donsker's
functional CLT, where entire random trajectories converge as curves) and to
\emph{non-asymptotic probability}. In modern high-dimensional statistics and machine
learning, an asymptotic guarantee that an error vanishes as $n \to \infty$ is often not
enough; researchers need explicit finite-sample bounds showing how many data points $n$
are required as a function of the dimension $d$, keeping convergence theory very much
alive.
\end{historybox}

\begin{paradoxbox}
\textbf{The St.\ Petersburg paradox and the limits of the LLN.} A casino offers a game:
flip a fair coin until the first head; if it appears on flip $k$, you win $2^k$ dollars.
The expected payout is $\sum_{k\ge1} 2^{-k}\cdot 2^k = \sum_{k\ge1}1 = \infty$. A
``rational'' player should pay any finite price to play. Yet nobody would pay a thousand
dollars per round. What broke?

The Law of Large Numbers broke -- or rather, its hypothesis did. The classical LLN
guarantees $\bar X_n\to\mu$ only when the mean $\mu$ is finite. Here $\mu=\infty$, so
the usual law gives no comfort. A subtler analysis (a weak law for these heavy-tailed
variables) shows that after $n$ rounds the \emph{typical} average winnings grow only
like $\log_2 n$, not like any finite constant -- the average drifts upward without
limit but agonizingly slowly, dominated by extremely rare colossal jackpots. So the
``infinite expectation'' is real but is delivered by events so rare that no finite
amount of play realizes it. This is the same ``rare giant'' phenomenon as
Example~\ref{ch22:ex-spike} in a gambling context: convergence statements
are only as good as their moment hypotheses, and a heavy tail can prevent convergence in
mean and probability alike.
\end{paradoxbox}

\begin{aestheticbox}
The four modes of convergence and their associated operational tools (continuous mapping,
Slutsky's theorem, and the delta method) form the foundational apparatus of asymptotic statistics.
When establishing the asymptotic normality of a studentized estimator,
$(\hat\theta-\theta)/\widehat{\mathrm{SE}} \dto\Normal(0,1)$, each piece plays a specific
mathematical role: the CLT provides convergence in distribution for the numerator,
the LLN provides convergence in probability to a constant for the estimated standard error
in the denominator, Slutsky's theorem combines them, and the delta method extends the conclusion
to smooth nonlinear transformations.
\end{aestheticbox}

\section{A worked synthesis}
\label{ch22:synthesis}

Here is the whole toolkit run once, end to end, on a single estimator, so the pieces
click together.

\begin{example}[Asymptotics of an estimated odds, start to finish]
\label{ch22:ex-odds}
You observe $n$ i.i.d.\ $\Bern(p)$ trials and want the asymptotic distribution of the
\emph{estimated odds} $\widehat{\text{odds}} = \hat p_n/(1-\hat p_n)$, where
$\hat p_n$ is the sample proportion. Watch each tool enter.

\textbf{Step 1 (LLN, $\pto$).} By the weak law, $\hat p_n\pto p$. So the estimator is
consistent: the odds estimate concentrates on the true odds $p/(1-p)$.

\textbf{Step 2 (CLT, $\dto$).} The standardized proportion is asymptotically normal:
\[
  \sqrt n\,(\hat p_n - p)\dto\Normal\!\big(0,\ p(1-p)\big).
\]

\textbf{Step 3 (delta method).} Let $g(x)=x/(1-x)$, so
$g'(x)=1/(1-x)^2$ and $g'(p)=1/(1-p)^2$. The delta method transports the normality:
\[
  \sqrt n\Big(\tfrac{\hat p_n}{1-\hat p_n}-\tfrac{p}{1-p}\Big)
  \dto \Normal\!\Big(0,\ \tfrac{1}{(1-p)^4}\cdot p(1-p)\Big)
  = \Normal\!\Big(0,\ \tfrac{p}{(1-p)^3}\Big).
\]

\textbf{Step 4 (Slutsky, for a usable interval).} The asymptotic variance
$p/(1-p)^3$ is unknown, but plugging in $\hat p_n$ gives a consistent estimate of it
(continuous mapping: the variance is a continuous function of $p$, and $\hat p_n\pto
p$). Dividing the centered, $\sqrt n$-scaled estimator by the square root of the
\emph{estimated} variance, Slutsky's quotient rule cancels the unknown and leaves a
standard normal limit. That yields the valid large-sample confidence interval
\[
  \frac{\hat p_n}{1-\hat p_n}\ \pm\ \frac{1.96}{\sqrt n}\,
  \sqrt{\frac{\hat p_n}{(1-\hat p_n)^3}}.
\]
Every one of the four modes and three tools made an appearance, in the natural order
they always appear: concentrate (LLN), find the fluctuation shape (CLT), push it
through your nonlinear quantity (delta method), and make it computable by swapping
unknowns for estimates (Slutsky).
\end{example}

\section{From static estimators to Markov chains}
\label{ch22:bridge}

Every sequence of random variables we analyzed in this chapter shared a quiet
common assumption: the underlying data points were independent, or the estimators
were static summaries of an i.i.d.\ sample drawn from an unchanging universe. The
target was always a fixed number $\mu$ or $\theta$, or a fixed universal distribution
like $\Normal(0,1)$ or $\chi^2_1$.

In the real world, systems have memory.
Tomorrow's weather depends on whether it is raining today. A financial market
tomorrow responds to where prices closed today. A molecule diffusing through a cell
takes its next step from its current location, not from scratch. When a sequence
$X_0, X_1, X_2, \dots$ unfolds through time with each step conditioned on the
present state, we enter the world of \emph{stochastic processes} -- random systems
evolving over time. The simplest and most widely used of these
processes is the \emph{Markov chain}, where the entire history of the past is
screened off by knowing where the system is right now.

When such a dynamic system runs for a long time, does it ever settle into balance?
The four modes of convergence from this chapter give us the exact vocabulary
to make that question precise:
\begin{itemize}
  \item \textbf{Convergence in distribution ($X_n \dto \pi$):} Does the probability
        distribution of where the system is located at time $n$ settle into a stable,
        time-invariant equilibrium -- a \emph{stationary distribution} $\pi$?
  \item \textbf{Almost sure convergence ($\frac{1}{n} \sum_{k=1}^n \ind_{\{X_k=i\}} \asto \pi_i$):}
        If you follow a \emph{single, unending simulation} through time, does the
        empirical fraction of time it spends visiting state $i$ lock onto $\pi_i$?
        This is the celebrated \emph{Ergodic Theorem} -- a Markovian version of the
        Strong Law of Large Numbers, asserting pathwise lock-in for dependent systems.
\end{itemize}

Without the four modes of convergence, asking whether an evolving system ``equilibrates''
would be hopelessly ambiguous. With them, we can cleanly separate the collective
balance of the population from the individual experience of a single path.
In Chapter~23, we take that step, moving from independent draws to processes whose next
step depends on where they stand right now.

\section{Exercises}
\label{ch22:exercises}

\begin{exercise}[Reading the definitions]
\label{ch22:exr-defs}
For each statement, name the strongest mode of convergence it asserts and explain in
one sentence why a weaker one would not capture it:
(a) ``With probability one, the running average lands on $\mu$ and stays.''
(b) ``At each large $n$, the estimator is within $0.01$ of the target with probability
at least $0.999$.''
(c) ``The mean-squared error tends to zero.''
(d) ``The histogram of the standardized statistic approaches a bell curve.''
\end{exercise}

\begin{exercise}[Markov to the rescue]
\label{ch22:exr-markov}
Suppose $\E(X_n - X)^2 = 1/n$. Prove directly from Markov's inequality that
$X_n\pto X$, and find an explicit $N$ such that $\P(\abs{X_n - X}\ge 0.1)\le 0.05$ for
all $n\ge N$.
\end{exercise}

\begin{exercise}[Build your own typewriter]
\label{ch22:exr-typewriter}
On $[0,1]$ with uniform probability, construct a sequence $X_n$ with $X_n\pto0$ but
such that $X_n(\omega)$ fails to converge for \emph{every} $\omega$. Verify both
claims. (You may reuse or modify the construction of
Section~\ref{ch22:typewriter}.)
\end{exercise}

\begin{exercise}[The rare giant, generalized]
\label{ch22:exr-giant}
On $[0,1]$ uniform, let $X_n = a_n\,\ind_{(0,1/n]}$. (a) For which sequences $a_n$ does
$X_n\pto0$? (b) For which does $X_n\to0$ in $L^1$? (c) In $L^2$? (d) Exhibit an $a_n$
for which $X_n\pto0$ but $X_n\not\to0$ in $L^1$, and another for which $X_n\to0$ in
$L^1$ but not $L^2$.
\end{exercise}

\begin{exercise}[Distribution does not imply probability]
\label{ch22:exr-d-not-p}
Let $Z\sim\Normal(0,1)$ and $X_n = (-1)^n Z$. Show $X_n\dto Z$ but $X_n\not\pto Z$.
Then show that if instead $X_n = Z + 1/n$, we do have $X_n\pto Z$. What single feature
of the first example breaks convergence in probability?
\end{exercise}

\begin{exercise}[Continuous mapping in action]
\label{ch22:exr-cmt}
Suppose $\sqrt n(\bar X_n - \mu)\dto\Normal(0,\sigma^2)$. Use the continuous mapping
theorem to find the limiting distribution of $n(\bar X_n - \mu)^2/\sigma^2$, and name
it.
\end{exercise}

\begin{exercise}[Slutsky, and a trap]
\label{ch22:exr-slutsky}
(a) Suppose $X_n\dto\Normal(0,1)$ and $Y_n\pto3$. Find the limiting distribution of
$Y_n X_n + Y_n$. (b) Now suppose instead $Y_n\dto\Normal(0,1)$ and you are told only
its marginal law. Explain carefully why you \emph{cannot} conclude
$X_n + Y_n\dto\Normal(0,2)$ in general, and describe what extra information would let
you.
\end{exercise}

\begin{exercise}[Delta method computations]
\label{ch22:exr-delta}
Suppose $\sqrt n(T_n-\theta)\dto\Normal(0,\sigma^2)$ with $\theta>0$. Find the
limiting distribution of $\sqrt n(\log T_n - \log\theta)$ and of
$\sqrt n(\sqrt{T_n}-\sqrt\theta)$. For the proportion estimator $\hat p_n$, derive the
asymptotic variance of the log-odds $\log(\hat p_n/(1-\hat p_n))$.
\end{exercise}

\begin{exercise}[Variance stabilization]
\label{ch22:exr-varstab}
For a $\Pois(\lambda)$ count estimator $\hat\lambda_n$ with
$\sqrt n(\hat\lambda_n-\lambda)\dto\Normal(0,\lambda)$, find a transform $g$ (a power
of $\hat\lambda_n$) whose delta-method variance does not depend on $\lambda$. (Hint:
try $g(x)=\sqrt x$ and compute $[g'(\lambda)]^2\lambda$.)
\end{exercise}

\begin{exercise}[Conceptual: why ``to a constant''?]
\label{ch22:exr-const}
Explain, with a concrete two-line counterexample, why Slutsky's theorem requires the
second sequence to converge to a constant rather than merely in distribution. (Hint:
consider $X_n=Z$ and $Y_n=-Z$ for a single fixed $Z\sim\Normal(0,1)$, versus
$Y_n=Z'$ an independent copy; look at $X_n+Y_n$.)
\end{exercise}

\begin{exercise}[Proof: $L^2 \Rightarrow L^1$]
\label{ch22:exr-l2-l1}
Using the Cauchy--Schwarz inequality (or Jensen's inequality applied to $x\mapsto
x^2$), prove that convergence in $L^2$ implies convergence in $L^1$. Then explain why
this, combined with Theorem~\ref{ch22:thm-lp-p}, means $L^2$ convergence sits high in
the hierarchy.
\end{exercise}

\begin{exercise}[Putting it together]
\label{ch22:exr-synthesis}
You estimate a rate $\hat r_n$ with $\sqrt n(\hat r_n - r)\dto\Normal(0,\tau^2)$ for a
known functional form $\tau^2=h(r)$. Outline, citing the specific theorem at each step,
how you would (a) argue $\hat r_n$ is consistent, (b) obtain the asymptotic
distribution of $g(\hat r_n)$ for a smooth $g$ with $g'(r)\ne0$, and (c) build an
asymptotic $95\%$ confidence interval for $g(r)$ that uses only observed quantities.
\end{exercise}
